\chapter{Conclusão}
\label{cap:conclusao}

O trabalho desenvolve um procedimento que separa métricas por imagem e por rosto, aceita estruturas de diretório arbitrárias e registra as condições da execução. Essa divisão evita repetir um mesmo escore CLIP conforme o número de faces detectadas. O pipeline também inclui meios de avaliar o detector e comparar as saídas automáticas com referências humanas e computacionais.

A conclusão empírica será finalizada depois da nova execução das 64 imagens e da consolidação do Ground Truth. Até essa etapa, não é possível afirmar se os \textit{prompts} inclusivos alteram as margens de diversidade racial percebida e de apresentação de gênero, nem se as evidências sustentam o Efeito Patchwork.

Os limites permanecem centrais para a leitura dos resultados. DeepFace e FairFace classificam aparências em faces sintéticas, enquanto CLIP mede alinhamento semântico com frases escolhidas pelo pesquisador. Nenhuma dessas ferramentas determina identidade demográfica. O \textit{framework} contribui ao registrar essas restrições junto com os resultados e ao tornar verificável o caminho entre imagem, medida e conclusão.
